\chapter{Improper Integrals}

In the previous chapter, we developed algebraic techniques---partial
fraction decomposition, trigonometric substitutions, and trigonometric integrals---to
find antiderivatives. Once you held an
antiderivative $F$, computing $\int_a^b f(x)\,\dd x$ felt like a solved puzzle:
evaluate $F(b) - F(a)$ and you were done.

Yet every integral you have met so far has been, in a quiet way, a finite object.
You integrate a \emph{bounded} function over a \emph{bounded} interval: $f$ does
not blow up, and the interval $[a,b]$ has two concrete endpoints you can point to
on the number line. Under those conditions the area under the curve is a single
finite number, the Riemann sums close in on it, and the Fundamental Theorem
hands you the answer through an antiderivative. Comfortable, complete, closed.

This chapter is about what happens when we break one or both of those
conditions on purpose. What is the area under a curve that runs off to the right
\emph{forever}? What is the area under a curve that shoots up to infinity as you
approach an endpoint? At first these questions sound either meaningless or
unavoidably infinite -- how could an infinitely wide region, or an infinitely tall
spike, enclose any finite amount of area?

\begin{focusbox}
The surprise -- and it genuinely surprised the people who first met it -- is that
sometimes an infinitely long or infinitely tall region has a perfectly finite
area, and sometimes it doesn't. Whether it does depends on \emph{how fast} the
curve shrinks or grows, in a way we can make completely precise. Learning to tell
the finite cases from the infinite ones is the central objective of this chapter.
\end{focusbox}

The device that lets us ask these questions rigorously is the same one that made
the derivative rigorous: the \textbf{limit}. We will never actually ``integrate to
infinity'' or ``integrate up to a vertical asymptote.'' Instead we will integrate
over an ordinary, finite, well-behaved region, and then \emph{watch what happens}
as we let the edge of that region march off to infinity or creep up to the bad
point. The integral over the infinite or unbounded region will be \emph{defined}
to be the limit of those ordinary integrals -- if the limit exists. That single
idea, applied carefully, organizes the entire chapter.

\chapterline

\section{Two ways an integral can be improper}
\label{ch25:two-ways}

Two things can go wrong, and naming them cleanly pays off because the cures
differ slightly.

\begin{definition}[Improper integral, informally]
\label{ch25:def-improper}
An integral $\int_a^b f(x)\,\dd x$ is called \textbf{improper} if either
\begin{enumerate}
  \item the interval of integration is infinite -- one or both limits are
  $\pm\infty$ (an improper integral \emph{of the first kind}); or
  \item the integrand $f$ is unbounded somewhere on the interval -- it has a
  vertical asymptote at an endpoint or at an interior point (an improper integral
  \emph{of the second kind}).
\end{enumerate}
In plain terms, either the interval of integration extends without bound, or the
integrand blows up somewhere along the interval.
An integral can be improper for both reasons at once.
\end{definition}

The ordinary Riemann integral simply has nothing to say about either case. Its
whole construction -- chop $[a,b]$ into finitely many pieces, build rectangles,
take a limit of the sums -- assumes a finite interval and a bounded function. If
the interval is $[1,\infty)$ there is no finite chopping to do; if $f$ blows up at
$x=0$, no finite collection of rectangles can pin down the area near the spike.
So we are not ``computing a Riemann integral'' in these cases. We are
\emph{defining a new kind of object} and choosing, deliberately, to define it as a
limit of ordinary integrals.

\begin{intuitionbox}
Think of it as photographing an open road or an endless coastline that extends past the
horizon. You cannot fit the whole stretch in one frame. So you take a frame that
stops at distance $t$, measure the area it contains, then slide the right edge of
the frame out -- $t = 10$, $t = 100$, $t = 10\,000$ -- and ask whether the measured
area is settling down toward a fixed number or growing without bound. The improper
integral is the number it settles toward.
\end{intuitionbox}

\begin{figure}[H]
\centering
\begin{tikzpicture}
\begin{axis}[
  width=12cm, height=5.2cm,
  domain=0.6:9.6, samples=120,
  axis lines=middle,
  xlabel={$x$}, ylabel={$y$},
  xtick={1,4}, xticklabels={$1$,$t$},
  ytick=\empty,
  ymin=0, ymax=1.25, xmin=0, xmax=10,
  enlargelimits=false, clip=false,
]
\addplot[curveblue, very thick, domain=0.7:9.6] {1/x^2};
\addplot[draw=none, fill=areafill, fill opacity=0.55, domain=1:4]
  {1/x^2} \closedcycle;
\draw[dashed, gray] (axis cs:4,0) -- (axis cs:4,0.0625);
\node[curveblue] at (axis cs:7.0,0.55) {$y = \dfrac{1}{x^2}$};
\node at (axis cs:2.5,0.72) {area $\displaystyle\int_1^{t}\tfrac{\dd x}{x^2}$};
\draw[->, thick] (axis cs:4.3,0.18) -- (axis cs:6.2,0.18)
  node[midway, above, font=\small] {let $t\to\infty$};
\end{axis}
\end{tikzpicture}
\caption{The frame that stops at $x=t$. We measure a finite area, then slide the
right edge $t$ outward and ask what number the area approaches.}
\label{ch25:fig-frame}
\end{figure}

\section{Infinite intervals: the first kind}
\label{ch25:first-kind}

Here is the definition, stated for the most common case of a single infinite
endpoint on the right.

\begin{definition}[Improper integral over $[a,\infty)$]
\label{ch25:def-first-kind}
Let $f$ be defined on $[a,\infty)$ and integrable on every finite subinterval
$[a,t]$. We define
\[
  \int_a^{\infty} f(x)\,\dd x
  \;=\; \lim_{t\to\infty} \int_a^{t} f(x)\,\dd x,
\]
\emph{provided this limit exists as a finite number.} In that case we say the
improper integral \textbf{converges} (the accumulated area settles down to a
single finite number), and its value is that limit. If the limit is
$\pm\infty$ or fails to exist, we say the integral \textbf{diverges} (the area
climbs without bound or oscillates without settling).
\end{definition}

Symmetrically, $\displaystyle\int_{-\infty}^{b} f(x)\,\dd x = \lim_{s\to-\infty}
\int_s^{b} f(x)\,\dd x$. The doubly-infinite case needs a little care, and we will
get to it.

\begin{happeningbox}
The symbol $\int_1^{\infty}$ looks like it has an endpoint called ``$\infty$'' that
you plug into an antiderivative. It does not. There is no number $\infty$ to plug
in, and the Riemann machinery never runs over an infinite interval. What is
\emph{really} happening is a two-step process hidden inside one piece of notation:
first do an ordinary integral up to a movable finite wall $t$, getting a
\emph{function of $t$}; then take the limit of that function as $t\to\infty$. The
``$\infty$'' in the upper limit is shorthand for ``take the limit.'' Every correct
calculation with these objects re-separates those two steps.
\end{happeningbox}

\subsection{The canonical example: two curves that look almost identical}
\label{ch25:canonical}

Run the definition on the cleanest possible pair. Consider
\[
  \int_1^{\infty} \frac{\dd x}{x^2}
  \qquad\text{and}\qquad
  \int_1^{\infty} \frac{\dd x}{x}.
\]
Both integrands are positive, both decrease toward zero as $x\to\infty$, and if
you sketch them they look like nearly the same gently-falling curve. Yet one
encloses a finite area and the other does not. Watch.

For the first, work to the finite wall $t$:
\[
  \int_1^{t} \frac{\dd x}{x^2}
  = \int_1^{t} x^{-2}\,\dd x
  = \eval{-\frac{1}{x}}{1}^{t}
  = -\frac{1}{t} + 1
  = 1 - \frac{1}{t}.
\]
Now slide the wall out. As $t\to\infty$, the term $1/t$ collapses to $0$, so
\[
  \int_1^{\infty} \frac{\dd x}{x^2}
  = \lim_{t\to\infty}\Bigl(1 - \tfrac{1}{t}\Bigr) = 1.
\]
The area under $1/x^2$ from $1$ all the way out to infinity is exactly $1$.
A region of \emph{infinite width} has area $1$.

Now the second, the same way:
\[
  \int_1^{t} \frac{\dd x}{x}
  = \eval{\ln x}{1}^{t}
  = \ln t - \ln 1 = \ln t.
\]
And $\ln t \to \infty$ as $t\to\infty$ -- slowly, but it never stops climbing. So
\[
  \int_1^{\infty} \frac{\dd x}{x} = \lim_{t\to\infty} \ln t = \infty
  \quad\text{(diverges).}
\]

\begin{intuitionbox}
What is the real difference? Both curves fall to zero, but $1/x^2$ falls
\emph{fast enough} that the leftover area out past $x=t$ shrinks to nothing,
while $1/x$ falls \emph{too slowly} -- it leaves behind a thin but stubbornly
non-negligible sliver of area at every stage, and those slivers add up to
infinity. The whole subject is the contest between ``the tail is getting thin''
and ``but there is infinitely much of it.'' Convergence means thinness wins.
\end{intuitionbox}

\begin{figure}[H]
\centering
\begin{tikzpicture}
\begin{axis}[
  width=12cm, height=5.6cm,
  domain=1:9.5, samples=160,
  axis lines=middle,
  xlabel={$x$}, ylabel={$y$},
  xtick={1}, ytick={1},
  ymin=0, ymax=1.15, xmin=0, xmax=10,
  enlargelimits=false, clip=false, legend pos=north east,
]
\addplot[curveblue, very thick] {1/x};
\addplot[tangentred, very thick] {1/x^2};
\legend{$1/x$ (area $=\infty$), $1/x^2$ (area $=1$)}
\addplot[draw=none, fill=areafill, fill opacity=0.35, domain=1:9.5]
  {1/x^2} \closedcycle;
\end{axis}
\end{tikzpicture}
\caption{Two curves that look nearly the same near the start. The faster-falling
$1/x^2$ (red, shaded) encloses finite area; the slower-falling $1/x$ (blue)
encloses infinite area. The eye cannot tell them apart out in the tail -- only the
arithmetic of the limit can.}
\label{ch25:fig-x-vs-x2}
\end{figure}

\begin{labbox}
Before reading on, compute $\displaystyle\int_1^{\infty} \frac{\dd x}{x^{3}}$ by
hand. Integrate to $t$, then let $t\to\infty$. (You should get the finite value
$\tfrac12$.) Then try $\displaystyle\int_1^{\infty} \frac{\dd x}{\sqrt{x}}$ and
watch it diverge. The point is to feel, in your fingers, exactly where the
$1/t^{\text{something}}$ either dies or survives in the limit.
\end{labbox}

\subsection{The $p$-test: the master pattern}
\label{ch25:p-test}

The two examples above are the $p=2$ and $p=1$ cases of one general family. Settle
the whole family at once, since it becomes the yardstick against which we measure
everything else.

\begin{theorem}[$p$-test at infinity]
\label{ch25:thm-p-infinity}
For a fixed real number $p$,
\[
  \int_1^{\infty} \frac{\dd x}{x^{p}}
  \quad\text{converges if and only if}\quad p > 1,
\]
and in the convergent case its value is $\dfrac{1}{p-1}$.
\end{theorem}

\begin{proof}
The only delicate part is that the antiderivative of $x^{-p}$ has a different form
when $p=1$, so we split the argument.

\emph{Case $p \ne 1$.} Then
\[
  \int_1^{t} x^{-p}\,\dd x
  = \eval{\frac{x^{1-p}}{1-p}}{1}^{t}
  = \frac{t^{1-p} - 1}{1-p}.
\]
Everything now rides on the exponent $1-p$ of $t$. If $p > 1$ then $1-p < 0$, so
$t^{1-p} = 1/t^{p-1} \to 0$ as $t\to\infty$, leaving
\[
  \lim_{t\to\infty}\frac{t^{1-p}-1}{1-p} = \frac{0-1}{1-p} = \frac{1}{p-1},
\]
a finite number: \emph{convergence}. If instead $p < 1$ then $1-p > 0$, so
$t^{1-p}\to\infty$ and the integral diverges.

\emph{Case $p = 1$.} Here the antiderivative is the logarithm, not a power:
\[
  \int_1^{t}\frac{\dd x}{x} = \ln t \to \infty,
\]
so the integral diverges. Collecting the cases: convergence happens exactly when
$p > 1$, with value $1/(p-1)$.
\end{proof}

\begin{whybox}
Why is $p=1$ the knife-edge -- the exact boundary between win and lose? Because
$1/x$ is the unique power whose antiderivative is logarithmic rather than a power.
A logarithm grows slower than \emph{any} positive power of $t$, yet it still grows
without bound. So $1/x$ sits precisely on the border: its tail is just barely too
fat to enclose finite area. Push the exponent up to $1.0000001$ and the
antiderivative becomes a (negative) power that converges; drop it to $0.9999999$
and you still diverge. The logarithm is the borderline case, and it falls on the
side of divergence. The threshold $p=1$ serves as the benchmark: convergence
questions over $[1,\infty)$ fundamentally reduce to whether the integrand decays
faster or slower than $1/x$.
\end{whybox}

\begin{trapbox}
``The integrand goes to zero, so the area must be finite.'' This is the single
most common mistake, and $\int_1^\infty \dd x/x$ is the clearest counterexample. The integrand
\emph{does} go to zero, and the integral still diverges. Going to zero is
\emph{necessary} for convergence over an infinite interval (loosely speaking), but
it is nowhere near sufficient -- the tail has to go to zero \emph{fast enough}.
Never conclude convergence just from $f(x)\to 0$.
\end{trapbox}

\subsection{Exponential tails and rapid decay}
\label{ch25:exp-tails}

While power functions like $1/x^2$ decay at a polynomial rate, exponential functions
decay far more rapidly and converge with room to spare. They appear frequently in
probability and physics, beginning with the standard decaying exponential.

\begin{example}[The decaying exponential]
\label{ch25:ex-exp}
Evaluate $\displaystyle\int_0^{\infty} e^{-x}\,\dd x$.
\[
  \int_0^{t} e^{-x}\,\dd x
  = \eval{-e^{-x}}{0}^{t}
  = -e^{-t} + e^{0}
  = 1 - e^{-t}.
\]
As $t\to\infty$, $e^{-t}\to 0$, so the integral converges to $1$. More generally,
for any rate $k>0$,
\[
  \int_0^{\infty} e^{-kx}\,\dd x
  = \lim_{t\to\infty}\eval{-\tfrac{1}{k}e^{-kx}}{0}^{t}
  = \frac{1}{k}.
\]
\end{example}

\begin{applicationbox}
This is the backbone of the \emph{exponential distribution} in probability, the
model for waiting times -- how long until a radioactive atom decays, until the next
customer arrives, until a component fails. The density is $f(x) = k e^{-kx}$ for
$x \ge 0$, and the calculation above says $\int_0^\infty k e^{-kx}\,\dd x = 1$:
total probability is one, exactly as a probability density must give. The whole
edifice of continuous probability rests on improper integrals converging to $1$.
The mean waiting time, $\int_0^\infty x\,k e^{-kx}\,\dd x = 1/k$, is another
improper integral -- you can do it with integration by parts and the same
$t\to\infty$ limit.
\end{applicationbox}

\subsection{Both ends infinite: split, don't straddle}
\label{ch25:both-ends}

What about $\int_{-\infty}^{\infty} f(x)\,\dd x$, infinite at \emph{both} ends? The
correct definition splits the line at any convenient finite point $c$ and demands
that \emph{each half} converge on its own.

\begin{definition}[Doubly infinite integral]
\label{ch25:def-doubly}
\[
  \int_{-\infty}^{\infty} f(x)\,\dd x
  = \int_{-\infty}^{c} f(x)\,\dd x + \int_{c}^{\infty} f(x)\,\dd x,
\]
and the left side converges if and only if \emph{both} integrals on the right
converge (the value is then independent of the split point $c$).
\end{definition}

\begin{trapbox}
You must \emph{not} compute $\int_{-\infty}^{\infty} f$ as a single symmetric limit
$\lim_{t\to\infty}\int_{-t}^{t} f$. That symmetric quantity is a different,
weaker object called the \emph{principal value} (specifically, the Cauchy
principal value---a single balanced limit where both boundaries are forced to
march outward in lockstep at the exact same rate: $\lim_{t\to\infty}\int_{-t}^t f(x)\,\dd x$).
The principal value can give a finite answer even when the true integral
diverges. Example: $\int_{-t}^{t} x\,\dd x = 0$ for every $t$, so the symmetric
limit is $0$. But $\int_0^\infty x\,\dd x = \infty$, so the genuine improper
integral $\int_{-\infty}^{\infty} x\,\dd x$ \emph{diverges}. The symmetric limit
was letting $+\infty$ and $-\infty$ cancel each other in lockstep, which is
cheating -- convergence is not allowed to depend on marching the two ends out at
the same speed. Always split into two independent halves.
\end{trapbox}

\begin{example}[A doubly-infinite winner]
\label{ch25:ex-arctan}
Evaluate $\displaystyle\int_{-\infty}^{\infty}\frac{\dd x}{1+x^2}$. Split at
$c=0$. The antiderivative is $\arctan x$, so
\[
  \int_0^{t}\frac{\dd x}{1+x^2} = \arctan t - \arctan 0 \to \frac{\pi}{2},
\]
and by the symmetry of the integrand the left half contributes another $\pi/2$.
Both halves converge, so the whole integral converges to
\[
  \int_{-\infty}^{\infty}\frac{\dd x}{1+x^2} = \frac{\pi}{2} + \frac{\pi}{2} = \pi.
\]
A bell-shaped curve stretching to infinity in both directions, total area exactly
$\pi$. The appearance of $\pi$ here is a cousin of why $\pi$ shows up in the
Gaussian and throughout probability.
\end{example}

\section{Unbounded integrands: the second kind}
\label{ch25:second-kind}

Now the other way an integral can be improper: the interval is finite, but the
\emph{function} runs off to infinity. Consider $\int_0^1 \dd x/\sqrt{x}$. The
interval $[0,1]$ is as tame as can be, but the integrand $1/\sqrt{x}\to\infty$ as
$x\to 0^+$. There is an infinitely tall spike pinned against the left edge. Does
it enclose finite area?

The cure is the same in spirit as before. We never integrate \emph{at} the bad
point. We back off to a wall a tiny distance $\eps$ away from it, integrate over
the safe region, and then let the wall creep in toward the singularity---a
\textbf{singularity} being any point where a function blows up to infinity or
ceases to be defined.

\begin{definition}[Singularity at an endpoint]
\label{ch25:def-second-kind}
Suppose $f$ is integrable on $[a+\eps,\,b]$ for every small $\eps>0$ but is
unbounded as $x\to a^+$. Then
\[
  \int_a^{b} f(x)\,\dd x
  = \lim_{\eps\to 0^+}\int_{a+\eps}^{b} f(x)\,\dd x,
\]
provided the limit exists (finite). Symmetrically, if the singularity is at the
\emph{right} endpoint $b$, back off from there:
$\displaystyle \int_a^{b} f = \lim_{\eps\to 0^+}\int_a^{b-\eps} f$.
\end{definition}

\begin{figure}[H]
\centering
\begin{tikzpicture}
\begin{axis}[
  width=11cm, height=5.6cm,
  domain=0.04:1, samples=160,
  axis lines=middle,
  xlabel={$x$}, ylabel={$y$},
  xtick={0.15,1}, xticklabels={$\eps$,$1$},
  ytick=\empty,
  ymin=0, ymax=5.2, xmin=-0.05, xmax=1.15,
  enlargelimits=false, clip=false,
]
\addplot[curveblue, very thick, domain=0.04:1] {1/sqrt(x)};
\addplot[draw=none, fill=areafill, fill opacity=0.55, domain=0.15:1]
  {1/sqrt(x)} \closedcycle;
\draw[dashed, gray] (axis cs:0.15,0) -- (axis cs:0.15,2.58);
\node[curveblue] at (axis cs:0.62,3.0) {$y = \dfrac{1}{\sqrt{x}}$};
\draw[->, thick] (axis cs:0.15,1.1) -- (axis cs:0.03,1.1)
  node[midway, above, font=\small] {$\eps\to 0^+$};
\end{axis}
\end{tikzpicture}
\caption{An unbounded integrand. The spike at $x=0$ is infinitely tall, so we cut
the area off at $x=\eps$, measure it, and slide the wall $\eps$ toward $0$. If the
measured area settles, the spike has finite area despite its infinite height.}
\label{ch25:fig-spike}
\end{figure}

\begin{example}[A tall spike with finite area]
\label{ch25:ex-sqrt}
Evaluate $\displaystyle\int_0^{1}\frac{\dd x}{\sqrt{x}}$. Back off to $\eps$:
\[
  \int_{\eps}^{1} x^{-1/2}\,\dd x
  = \eval{2\sqrt{x}}{\eps}^{1}
  = 2\sqrt{1} - 2\sqrt{\eps}
  = 2 - 2\sqrt{\eps}.
\]
As $\eps\to 0^+$, $\sqrt{\eps}\to 0$, so the integral converges to $2$. The spike is
infinitely tall, and it encloses area $2$.
\end{example}

\begin{example}[A tall spike with infinite area]
\label{ch25:ex-recip}
Evaluate $\displaystyle\int_0^{1}\frac{\dd x}{x}$. The same procedure:
\[
  \int_{\eps}^{1}\frac{\dd x}{x} = \eval{\ln x}{\eps}^{1}
  = \ln 1 - \ln\eps = -\ln\eps.
\]
As $\eps\to 0^+$, $\ln\eps\to -\infty$, so $-\ln\eps\to+\infty$: \emph{diverges}.
Again the logarithm sits on the knife-edge, again it falls on the side of
divergence.
\end{example}

\subsection{A $p$-test for spikes -- and a beautiful inversion}
\label{ch25:p-test-zero}

There is a $p$-test here too, and comparing it to the one at infinity reveals one
of the prettiest little symmetries in calculus.

\begin{theorem}[$p$-test at a singularity]
\label{ch25:thm-p-zero}
For a fixed real $p$,
\[
  \int_0^{1} \frac{\dd x}{x^{p}}
  \quad\text{converges if and only if}\quad p < 1,
\]
with value $\dfrac{1}{1-p}$ in the convergent case.
\end{theorem}

\begin{proof}
For $p\ne 1$,
\[
  \int_{\eps}^{1} x^{-p}\,\dd x
  = \eval{\frac{x^{1-p}}{1-p}}{\eps}^{1}
  = \frac{1 - \eps^{\,1-p}}{1-p}.
\]
If $p < 1$ then $1-p > 0$ and $\eps^{\,1-p}\to 0$, giving the finite limit
$1/(1-p)$. If $p > 1$ then $1-p < 0$ and $\eps^{\,1-p} = 1/\eps^{\,p-1}\to\infty$,
so it diverges. For $p=1$ we got $-\ln\eps\to\infty$ above. Hence convergence
exactly when $p<1$.
\end{proof}

\begin{aestheticbox}
Lay the two $p$-tests side by side and admire the reversal:
\[
  \int_1^{\infty}\frac{\dd x}{x^p}\ \text{converges} \iff p > 1,
  \qquad
  \int_0^{1}\frac{\dd x}{x^p}\ \text{converges} \iff p < 1.
\]
The inequality literally flips. And it \emph{must}, because the two integrals ask
opposite questions about the same exponent. Far out at infinity, you want the
curve to fall \emph{quickly} (large $p$) so the long tail is thin. Near zero, you
want the curve to rise \emph{slowly} (small $p$) so the tall spike is narrow.
A large $p$ makes a thin tail but a fat spike; a small $p$ makes a thin spike but a
fat tail. You cannot have both -- which is exactly why $\int_0^\infty \dd x/x^p$
\emph{diverges for every $p$}: whatever value of $p$ tames one end ruins the
other. The borderline curve $1/x$ ($p=1$) diverges at both ends at once.
\end{aestheticbox}

\begin{recoverybox}
\textbf{Mental check:} Without writing anything down, which of these two integrals converges:
$\displaystyle\int_0^1 \frac{\dd x}{x^{0.99}}$ or $\displaystyle\int_0^1 \frac{\dd x}{x^{1.01}}$?
And which one converges if the interval is flipped to $[1,\infty)$?

\emph{Answer:} Near zero, the smaller power $p = 0.99 < 1$ climbs slowly enough that the vertical spike is thin, so $\int_0^1 \dd x/x^{0.99}$ converges ($= 1/(1-0.99) = 100$). The power $p = 1.01 > 1$ blows up too aggressively and diverges. Out on the infinite tail $[1,\infty)$, the roles reverse: $p = 1.01 > 1$ falls fast enough to enclose finite area ($= 1/(1.01-1) = 100$), while $p = 0.99$ diverges. The boundary $p = 1$ is the logarithm, which diverges in both directions.
\end{recoverybox}

\begin{figure}[H]
\centering
\begin{tikzpicture}
% Left panel: Spike on (0, 1]
\begin{axis}[
  name=spike,
  width=7.2cm, height=5.6cm,
  axis lines=left,
  xlabel={$x$}, ylabel={$y$},
  title style={font=\small},
  title={\textbf{(a)} Singularity: $\int_0^1 x^{-p}\,\dd x$ ($p<1$ converges)},
  xtick={0.2, 1.0}, xticklabels={$\eps$, $1$},
  ytick=\empty,
  xmin=0, xmax=1.2, ymin=0, ymax=4.2,
  domain=0.06:1.0, samples=80,
  clip=false,
]
% Shaded convergent spike p=0.5
\addplot[draw=none, fill=targetgreen, opacity=0.25, domain=0.06:1.0] {x^(-0.5)} \closedcycle;
\addplot[targetgreen!80!black, very thick, domain=0.06:1.0] {x^(-0.5)};
% Divergent spike p=2
\addplot[trapred, thick, dashed, domain=0.48:1.0] {x^(-2)};

\draw[dashed, gray] (axis cs:0.2,0) -- (axis cs:0.2,2.23);
\node[targetgreen!80!black, font=\footnotesize, anchor=north west] at (axis cs:0.25,1.0) {$p=\tfrac12$ (area $=2$)};
\node[trapred!80!black, font=\footnotesize, anchor=south west] at (axis cs:0.65,3.0) {$p=2$ ($\infty$)};
\end{axis}

% Right panel: Tail on [1, \infty)
\begin{axis}[
  name=tail,
  at={($(spike.east)+(1.6cm,0)$)},
  anchor=west,
  width=7.2cm, height=5.6cm,
  axis lines=left,
  xlabel={$x$}, ylabel={},
  title style={font=\small},
  title={\textbf{(b)} Infinite tail: $\int_1^\infty x^{-p}\,\dd x$ ($p>1$ converges)},
  xtick={1.0, 4.0}, xticklabels={$1$, $t$},
  ytick=\empty,
  xmin=0.8, xmax=6.5, ymin=0, ymax=1.3,
  domain=1.0:6.2, samples=80,
  clip=false,
]
% Shaded convergent tail p=2
\addplot[draw=none, fill=sourceblue, opacity=0.25, domain=1.0:6.0] {x^(-2)} \closedcycle;
\addplot[sourceblue!80!black, very thick, domain=1.0:6.0] {x^(-2)};
% Divergent tail p=0.5
\addplot[trapred, thick, dashed, domain=1.0:6.0] {x^(-0.5)};

\draw[dashed, gray] (axis cs:4.0,0) -- (axis cs:4.0,0.0625);
\node[sourceblue!80!black, font=\footnotesize, anchor=west] at (axis cs:1.8,0.42) {$p=2$ (area $=1$)};
\node[trapred!80!black, font=\footnotesize, anchor=south west] at (axis cs:2.5,0.85) {$p=\tfrac12$ ($\infty$)};
\end{axis}

\end{tikzpicture}
\caption{The trade-off behind the flipped $p$-tests. \textbf{(a)} Near a vertical singularity on $(0,1]$, small exponents $p < 1$ converge (green shaded, area $= 2$) because the spike is narrow, whereas large exponents $p \ge 1$ diverge (dashed red). \textbf{(b)} Out on the infinite tail $[1,\infty)$, large exponents $p > 1$ converge (blue shaded, area $= 1$) because the tail decays rapidly, whereas small exponents $p \le 1$ diverge (dashed red). No single $p$ can win at both ends -- which is why $\int_0^\infty \dd x/x^p$ diverges for every $p$.}
\label{ch25:fig-tradeoff}
\end{figure}

\section{When you cannot find an antiderivative: comparison}
\label{ch25:comparison}

So far every example handed us a clean antiderivative. But the questions we most
want to answer often involve integrands like $e^{-x^2}$ or
$\dfrac{1}{\sqrt{x^3+1}}$ that have \emph{no elementary antiderivative} at all. We
cannot evaluate $\int_1^\infty e^{-x^2}\,\dd x$ by finding $F$ and taking a limit --
there is no $F$ to find. Yet we can still answer the only question that matters
for many purposes: does it \emph{converge}?

The idea is the oldest trick in analysis: if you cannot measure a quantity, trap
it between two quantities you \emph{can} measure.

\begin{intuitionbox}
If your messy curve sits \emph{underneath} a curve whose area you know to be
finite, then your curve -- having even less area -- must also have finite area. And
if your messy curve sits \emph{above} a curve whose area is infinite, then yours,
having even more, must be infinite too. You never compute your own area; you
sandwich it. This is the same move as the Squeeze Theorem for limits, now applied
to whole accumulated areas.
\end{intuitionbox}

\begin{theorem}[Comparison Test for integrals]
\label{ch25:thm-comparison}
Suppose $0 \le f(x) \le g(x)$ for all $x \ge a$, and both are integrable on every
$[a,t]$. Then:
\begin{enumerate}
  \item If $\displaystyle\int_a^{\infty} g$ converges, so does
  $\displaystyle\int_a^{\infty} f$ (the smaller one is also finite).
  \item If $\displaystyle\int_a^{\infty} f$ diverges, so does
  $\displaystyle\int_a^{\infty} g$ (the bigger one is also infinite).
\end{enumerate}
The same statements hold for second-kind (unbounded-integrand) integrals.
\end{theorem}

\begin{proof}
The assumption $f \ge 0$ is essential for monotonicity. Define the accumulated area
function
\[
  A(t) = \int_a^{t} f(x)\,\dd x.
\]
Because $f \ge 0$, adding more interval only adds (non-negative) area, so $A$ is a
\emph{nondecreasing} function of $t$. A nondecreasing function as $t\to\infty$ has
only two possible fates: it climbs to a finite limit, or it runs off to $+\infty$.
(There is no oscillation to worry about -- it cannot come back down.) This is the
Monotone Convergence principle (the guarantee that any quantity that only ever
grows and cannot cross a fixed ceiling must settle into a finite limit), and it
is the engine of the whole proof.

For part (1): since $f \le g$, comparing areas over the same interval gives
\[
  A(t) = \int_a^{t} f \;\le\; \int_a^{t} g \;\le\; \int_a^{\infty} g =: M,
\]
where $M$ is finite by hypothesis. So $A$ is nondecreasing \emph{and bounded above}
by $M$. A nondecreasing function bounded above must converge to a finite limit
(it cannot escape past its ceiling). Hence $\int_a^\infty f$ converges.

Part (2) is the contrapositive of part (1): if $\int_a^\infty f$ diverged to
$+\infty$ while $\int_a^\infty g$ converged, part (1) would force $\int_a^\infty f$
to converge -- a contradiction. So if $\int f$ diverges, $\int g$ must diverge too.
\end{proof}

\begin{whybox}
Why does the comparison test \emph{need} $f$ and $g$ to be non-negative? Because
the whole argument rests on $A(t)$ being monotone -- on the fact that the area
``so far'' can only grow as we extend the interval. If $f$ were allowed to dip
negative, $A(t)$ could rise and fall, and ``bounded above'' would no longer force
convergence (a bounded oscillation need not settle anywhere). Sign-control is what
converts ``bounded'' into ``convergent.'' For integrands that change sign you need
the stronger notion of \emph{absolute convergence}, coming up at the end of the
chapter.
\end{whybox}

\begin{example}[Trapping a curve with no antiderivative]
\label{ch25:ex-gaussian-tail}
Show $\displaystyle\int_1^{\infty} e^{-x^2}\,\dd x$ converges. There is no
elementary antiderivative, so we cannot evaluate it -- but for $x \ge 1$ we have
$x^2 \ge x$, hence $e^{-x^2} \le e^{-x}$. And we computed
$\int_1^\infty e^{-x}\,\dd x = e^{-1}$, which is finite. By part (1) of the
Comparison Test, our integral is squeezed under a finite area and therefore
converges. We have proven it is finite without ever finding its value.
\end{example}

\begin{figure}[H]
\centering
\begin{tikzpicture}
\begin{axis}[
  width=12cm, height=5.4cm,
  domain=1:5, samples=120,
  axis lines=middle,
  xlabel={$x$}, ylabel={$y$},
  xtick={1}, ytick=\empty,
  ymin=0, ymax=0.42, xmin=0, xmax=5.4,
  enlargelimits=false, clip=false, legend pos=north east,
]
\addplot[curveblue, very thick] {exp(-x)};
\addplot[tangentred, very thick] {exp(-x^2)};
\addplot[draw=none, fill=areafill, fill opacity=0.45, domain=1:5]
  {exp(-x^2)} \closedcycle;
\legend{$e^{-x}$ (finite area), $e^{-x^2}$ (squeezed below)}
\end{axis}
\end{tikzpicture}
\caption{The comparison test in one picture. For $x\ge 1$ the curve $e^{-x^2}$
(red, shaded) lies entirely under $e^{-x}$ (blue), whose tail area is finite. A
curve trapped beneath a finite area must itself have finite area -- so
$\int_1^\infty e^{-x^2}\,\dd x$ converges, even though we never compute it.}
\label{ch25:fig-comparison}
\end{figure}

\begin{example}[Using divergence to prove divergence]
\label{ch25:ex-divcomp}
Show $\displaystyle\int_1^{\infty}\frac{\dd x}{\sqrt{x^2+x}}$ diverges. For
$x\ge 1$ we have $x^2 + x \le x^2 + x^2 = 2x^2$, so
\[
  \frac{1}{\sqrt{x^2+x}} \;\ge\; \frac{1}{\sqrt{2x^2}} = \frac{1}{\sqrt{2}\,x}.
\]
But $\int_1^\infty \frac{1}{\sqrt 2\,x}\,\dd x = \frac{1}{\sqrt2}\int_1^\infty
\frac{\dd x}{x}$ diverges (it is the $p=1$ borderline). Our integrand sits
\emph{above} a divergent one, so by part (2) it diverges too.
\end{example}

\begin{trapbox}
Comparison only works in the helpful direction. Bounding $f$ \emph{above} by a
\emph{divergent} $g$ tells you nothing -- a small thing under a big infinite thing
can be finite or infinite. Likewise bounding $f$ \emph{below} by a \emph{convergent}
$g$ tells you nothing. The useful pairings are: ``smaller than convergent
$\Rightarrow$ convergent'' and ``bigger than divergent $\Rightarrow$ divergent.''
Get the direction backwards and you have proven nothing. When in doubt, ask: is my
known curve acting as a finite ceiling, or an infinite floor?
\end{trapbox}

\subsection{Limit comparison: comparing by behavior, not by inequality}
\label{ch25:limit-comparison}

Building an explicit inequality $f \le g$ can be fiddly -- you have to nudge the
algebra into the right direction. Often you can \emph{see} that two integrands
behave the same way far out (both ``act like'' $1/x^2$, say) even though neither is
cleanly below the other. The limit-comparison test makes that intuition rigorous.

\begin{theorem}[Limit Comparison Test]
\label{ch25:thm-limit-comparison}
Let $f, g > 0$ on $[a,\infty)$ and suppose
\[
  \lim_{x\to\infty}\frac{f(x)}{g(x)} = L
  \quad\text{with } 0 < L < \infty.
\]
Then $\int_a^\infty f$ and $\int_a^\infty g$ either both converge or both diverge:
they share the same fate.
\end{theorem}

The hypothesis $0 < L < \infty$ says $f$ and $g$ are, far out, the same size up to
a bounded constant factor -- and a bounded constant factor cannot turn a finite
area infinite or vice versa. (We sketch the proof in an exercise; it is a clean
application of the plain comparison test, bounding $f$ between $\frac{L}{2}g$ and
$2L\,g$ for large $x$.)

\begin{example}[Limit comparison in action]
\label{ch25:ex-limcomp}
Does $\displaystyle\int_1^{\infty}\frac{3x^2 - x + 5}{2x^4 + 7}\,\dd x$ converge?
Forget the constants and look at the dominant powers: for large $x$ the integrand
behaves like $\dfrac{3x^2}{2x^4} = \dfrac{3}{2x^2}$. So compare against
$g(x) = 1/x^2$:
\[
  \lim_{x\to\infty}\frac{(3x^2-x+5)/(2x^4+7)}{1/x^2}
  = \lim_{x\to\infty}\frac{(3x^2-x+5)\,x^2}{2x^4+7}
  = \lim_{x\to\infty}\frac{3x^4 - x^3 + 5x^2}{2x^4 + 7}
  = \frac{3}{2}.
\]
The limit is finite and positive, so our integral shares the fate of
$\int_1^\infty \dd x/x^2$, which \emph{converges} ($p=2 > 1$). And we never had to
integrate the ugly rational function.
\end{example}

\begin{focusbox}
This provides a direct criterion for rational functions. For an integrand $P(x)/Q(x)$,
compare degrees: if the denominator degree exceeds the numerator degree by more than one
($\deg Q - \deg P > 1$), the integrand decays faster than $1/x$ and the integral over
$[1,\infty)$ converges. If $\deg Q - \deg P \le 1$, the integral diverges. The difference
$\deg Q - \deg P$ plays the role of the exponent $p$ in the $p$-test.
\end{focusbox}

\begin{recoverybox}
\textbf{Mental check:} Suppose $0 \le f(x) \le \dfrac{1}{x^3}$ for all $x \ge 1$. Does $\displaystyle\int_1^\infty f(x)\,\dd x$ converge? What if you only knew $f(x) \ge \dfrac{1}{x^3}$?

\emph{Answer:} In the first case, $f$ is trapped under $1/x^3$, whose integral converges ($p = 3 > 1$). Having a finite ceiling forces $\int_1^\infty f$ to converge. In the second case, knowing $f$ sits above a convergent curve tells you nothing: $f(x)$ could be $1/x^4$ (convergent) or $1/x$ (divergent). A finite floor cannot stop a curve from escaping to infinity.
\end{recoverybox}

\section{The famous surprises}
\label{ch25:surprises}

We have already met one mild surprise -- finite area under an infinitely long
curve. The historical excitement around improper integrals came from a few results
that genuinely offended common sense. They are worth meeting both for the delight
and because they sharpen what ``finite area'' really does and does not promise.

\subsection{Finite area, infinite length}
\label{ch25:area-vs-length}

\begin{intuitionbox}
Take the region under $y = 1/x^2$ from $1$ to $\infty$: area $= 1$, finite. Now ask
about the \emph{arc length} of a curve that runs to infinity horizontally -- how
long is the wiggly line, not the area under it. The arc length formula gives
$\int_1^\infty \sqrt{1 + (y')^2}\,\dd x$. Since $\sqrt{1+(y')^2} \ge 1$ for
\emph{every} curve, the arc length is at least $\int_1^\infty 1\,\dd x = \infty$.
\emph{Any} curve that runs to infinity horizontally has infinite length, no matter
how flat it gets. So length and area are genuinely different beasts: one can be
infinite while the other is finite, on the very same picture.
\end{intuitionbox}

\subsection{Gabriel's Horn: finite volume, infinite surface}
\label{ch25:gabriels-horn}

A classic geometric demonstration discovered by Evangelista Torricelli around 1641 -- before
calculus had its modern notation -- illustrates this distinction directly. Rotate the curve $y = 1/x$ for $x \ge 1$
about the $x$-axis to produce the surface of revolution known as
\emph{Gabriel's Horn} (also called Torricelli's trumpet).

\begin{figure}[H]
\centering
\begin{tikzpicture}
\begin{axis}[
  width=12cm, height=4.6cm,
  axis lines=middle,
  xlabel={$x$},
  xtick={1}, ytick=\empty,
  ymin=-1.2, ymax=1.2, xmin=0, xmax=10.2,
  enlargelimits=false, clip=false,
  domain=1:9.2, samples=120,
]
\addplot[curveblue, very thick] {1/x};
\addplot[curveblue, very thick] {-1/x};
\draw[curveblue, thick] (axis cs:1,0) ellipse [x radius=0.12, y radius=1];
\node[curveblue] at (axis cs:5,0.55) {$y = 1/x$ spun about the $x$-axis};
\end{axis}
\end{tikzpicture}
\caption{Gabriel's Horn: the curve $y=1/x$ revolved about the $x$-axis, narrowing
forever as $x\to\infty$. Finite volume, infinite surface area.}
\label{ch25:fig-horn}
\end{figure}

Two improper integrals describe the horn. Using the disk and surface-of-revolution
formulas (which you will derive in the applications chapter; we just use them here):

\emph{Volume} (stacking up thin disks of radius $1/x$):
\[
  V = \int_1^{\infty} \pi\,\Bigl(\frac{1}{x}\Bigr)^2 \dd x
    = \pi \int_1^{\infty}\frac{\dd x}{x^2}
    = \pi \cdot 1 = \pi.
\]
A finite volume. You could fill the entire infinitely-long horn with exactly $\pi$
cubic units of paint.

\emph{Surface area} (wrapping a thin band of circumference $2\pi/x$ around each
slice):
\[
  S = \int_1^{\infty} 2\pi\,\frac{1}{x}\sqrt{1 + \Bigl(\frac{-1}{x^2}\Bigr)^2}\,\dd x
    > \int_1^{\infty} \frac{2\pi}{x}\,\dd x = 2\pi \int_1^{\infty}\frac{\dd x}{x}
    = \infty.
\]
Here the surface integrand is at least $2\pi/x$ (because the square root is at
least $1$), and $\int_1^\infty \dd x/x$ diverges -- our old friend the $p=1$
borderline. So the surface area is \emph{infinite}.

\begin{aestheticbox}
The apparent paradox arises from contrasting volume with surface area. Gabriel's horn
encloses a finite volume of $\pi$, yet has infinite surface area. In physical terms,
filling the interior requires a finite volume of paint, whereas coating the surface
with a layer of uniform thickness $\delta > 0$ requires infinite volume (since the volume
of a coating of thickness $\delta$ is at least $\delta \cdot S = \infty$). The geometric
resolution is that volume and surface area are independent dimensional quantities:
\emph{a finite three-dimensional volume can be bounded by an infinite two-dimensional surface},
just as a finite area can be bounded by a curve of infinite length.
\end{aestheticbox}

\begin{historybox}
Evangelista Torricelli, a student of Galileo Galilei, analyzed this solid in 1641.
He calculated the volume and surface area of the infinitely long horn obtained by
rotating the hyperbola $y=1/x$ about its asymptote.

Torricelli worked before the development of calculus, using Cavalieri's \emph{method
of indivisibles}---comparing the solid to concentric cylindrical shells of known volume.
The result surprised contemporary scholars. Classical geometry, following Aristotle,
avoided completed infinities, and the idea that a solid of infinite length could enclose
a finite volume seemed counterintuitive. Thomas Hobbes argued against the result, while
René Descartes and Marin Mersenne debated its implications in correspondence.

Torricelli's calculation showed that geometric dimensions behave independently: infinite
length or surface area does not imply infinite volume.
\end{historybox}

\begin{connectionbox}
Gabriel's horn is the continuous geometric counterpart of a fundamental fact about infinite
series: $\sum 1/n^2$ converges, whereas the harmonic series $\sum 1/n$ diverges. The
Integral Test, developed in the series chapters, establishes that an improper integral
$\int_1^\infty f(x)\,\dd x$ and the corresponding infinite series $\sum f(n)$ for a
positive decreasing function $f$ share the exact same convergence behavior. Improper
integrals and infinite series both address the accumulation of infinitely many shrinking
terms. The $p$-test for improper integrals directly anticipates the $p$-series test.
\end{connectionbox}

\section{Hidden singularities and other traps}
\label{ch25:hidden}

The most dangerous improper integrals are the ones that do not announce themselves.
A finite interval and a formula that looks bounded can still hide a singularity in
the middle, and if you integrate straight through it with the Fundamental Theorem
you can get a confident, plausible, and completely wrong answer.

\begin{trapbox}
Consider the innocent-looking
\[
  \int_{-1}^{1} \frac{\dd x}{x^2}.
\]
A careless application of the Fundamental Theorem gives
$\eval{-1/x}{-1}^{1} = (-1) - (1) = -2$. This is \emph{nonsense}: the integrand
$1/x^2$ is positive everywhere it is defined, so any genuine ``area'' must be
positive, never $-2$. What went wrong? The integrand blows up at $x=0$, an
\emph{interior} point of $[-1,1]$, so the integral is improper and the
Fundamental Theorem -- which assumes a continuous integrand -- does not apply. The
formula $-1/x$ is not a valid antiderivative \emph{across} the discontinuity. Doing
it correctly, you split at $0$ and find each half is $\int_0^1 \dd x/x^2 = \infty$,
so the integral diverges. Always scan the \emph{whole} interval, interior included,
for points where the integrand is undefined before you reach for an antiderivative.
\end{trapbox}

\begin{figure}[H]
\centering
\begin{tikzpicture}
\begin{axis}[
  width=11cm, height=5.2cm,
  axis lines=middle,
  xlabel={$x$}, ylabel={$y$},
  xtick={-1, 1}, xticklabels={$-1$, $1$},
  ytick=\empty,
  ymin=0, ymax=6.5, xmin=-1.4, xmax=1.4,
  enlargelimits=false, clip=false,
]
\addplot[curveblue, very thick, domain=-1.3:-0.39, samples=80] {1/x^2};
\addplot[curveblue, very thick, domain=0.39:1.3, samples=80] {1/x^2};
\addplot[draw=none, fill=areafill, fill opacity=0.45, domain=-1:-0.41, samples=60] {1/x^2} \closedcycle;
\addplot[draw=none, fill=areafill, fill opacity=0.45, domain=0.41:1, samples=60] {1/x^2} \closedcycle;
\draw[dashed, trapred, thick] (axis cs:0,0) -- (axis cs:0,6.0);
\node[trapred!85!black, font=\small, anchor=west] at (axis cs:0.06,5.4) {singularity at $x=0$};
\node[curveblue, font=\small] at (axis cs:0.8,3.2) {$y = \dfrac{1}{x^2}$};
\end{axis}
\end{tikzpicture}
\caption{The hidden interior singularity trap: $\int_{-1}^1 \dd x/x^2$. The curve $y=1/x^2$ is strictly positive everywhere, shooting up to $+\infty$ on both sides of $x=0$. Applying the Fundamental Theorem blindly across the tear yields $-2$---an absurd negative area for a curve that never drops below the $x$-axis. The correct value is $\infty + \infty = \infty$ (diverges).}
\label{ch25:fig-hidden-singularity}
\end{figure}

\begin{definition}[Interior singularity]
\label{ch25:def-interior}
If $f$ is unbounded near an interior point $c$ of $[a,b]$, define
\[
  \int_a^{b} f = \int_a^{c} f + \int_c^{b} f,
\]
and the whole integral converges only if \emph{both} pieces converge as improper
integrals in their own right. As with two-sided infinity, you may not let the two
sides cancel across the singularity.
\end{definition}

\begin{example}[Careful handling of a spike]
\label{ch25:ex-interior}
Evaluate $\displaystyle\int_0^{8}\frac{\dd x}{x^{1/3}}$. The integrand blows up at
$x=0$ (an endpoint here, but the moral is the same). Back off and take the limit:
\[
  \int_{\eps}^{8} x^{-1/3}\,\dd x
  = \eval{\tfrac{3}{2}x^{2/3}}{\eps}^{8}
  = \tfrac{3}{2}\bigl(8^{2/3} - \eps^{2/3}\bigr)
  = \tfrac{3}{2}\bigl(4 - \eps^{2/3}\bigr).
\]
As $\eps\to 0^+$, $\eps^{2/3}\to 0$, so the value is $\tfrac{3}{2}\cdot 4 = 6$.
Here $p = 1/3 < 1$, so the $p$-test for singularities promised convergence, and
indeed it converges to $6$.
\end{example}

\begin{walkbox}
Take a walk and turn this over: the integrand $1/x$ diverges at the singularity
$x=0$ (over $[0,1]$) \emph{and} diverges at infinity (over $[1,\infty)$). The
exponent that would fix one end ($p$ large near $\infty$, $p$ small near $0$) is
exactly the wrong fix for the other. Convince yourself, without writing anything,
that there is \emph{no} power $1/x^p$ whose integral over the whole ray
$(0,\infty)$ is finite. What does that tell you about trying to assign a finite
``total area under $1/x^p$''? And what is special about $p=1$ that makes it fail
symmetrically at both ends? This is the whole chapter compressed into one curve.
\end{walkbox}

\section{When the sign keeps changing: absolute convergence}
\label{ch25:absolute}

Everything in the comparison section assumed $f \ge 0$, so that the area-so-far
$A(t)$ could only climb. What about an integrand that oscillates -- positive, then
negative, then positive -- like $\dfrac{\sin x}{x}$ out to infinity? Now $A(t)$
rises and falls, and ``does the area settle down?'' becomes a subtler question of
\emph{cancellation}, not just of size.

\begin{definition}[Absolute convergence]
\label{ch25:def-absolute}
The integral $\int_a^\infty f$ \textbf{converges absolutely} if
$\int_a^\infty \abs{f}$ converges. A basic theorem of analysis says: \emph{if
$\int_a^\infty \abs{f}$ converges, then $\int_a^\infty f$ converges too.} Absolute
convergence is the safe, robust kind -- and what makes it so useful in practice is that it is checkable by the
comparison test, because $\abs{f} \ge 0$.
\end{definition}

\begin{intuitionbox}
The reason absolute convergence implies ordinary convergence: if even the
\emph{total amount of up-and-down motion} (the area of $\abs{f}$, counting all
wiggles as positive) is finite, then the actual signed area, which enjoys helpful
cancellation on top of that, certainly cannot run off to infinity. Bounding the
total size bounds the signed total. The converse fails: an integral can converge
only \emph{because} of cancellation, with the absolute version diverging. That
fragile kind is called \emph{conditional} convergence, and $\int_1^\infty
\frac{\sin x}{x}\,\dd x$ is the classic example -- it converges, but $\int_1^\infty
\frac{\abs{\sin x}}{x}\,\dd x$ diverges, because the bumps of $\abs{\sin x}/x$ are,
on average, a constant times $1/x$.
\end{intuitionbox}

\begin{example}[Absolute convergence via comparison]
\label{ch25:ex-abs}
Show $\displaystyle\int_1^{\infty}\frac{\sin x}{x^2}\,\dd x$ converges. We cannot
sign-control $\sin x$, but we can bound its size: $\abs{\sin x}\le 1$, so
\[
  \abs{\frac{\sin x}{x^2}} \le \frac{1}{x^2}.
\]
Since $\int_1^\infty \dd x/x^2$ converges, the comparison test makes
$\int_1^\infty \abs{\sin x / x^2}$ converge -- that is, our integral converges
\emph{absolutely}, and hence converges. The oscillation never even had a chance to
cause trouble, because the envelope $1/x^2$ already had finite area.
\end{example}

\begin{figure}[H]
\centering
\begin{tikzpicture}
\begin{axis}[
  width=12cm, height=5.2cm,
  domain=1:10, samples=200,
  axis lines=middle,
  xlabel={$x$}, ylabel={$y$},
  xtick={1, 3.1416, 6.2832, 9.4248},
  xticklabels={$1$, $\pi$, $2\pi$, $3\pi$},
  ytick={-1, 1},
  ymin=-1.15, ymax=1.15, xmin=0, xmax=10.5,
  enlargelimits=false, clip=false,
  legend style={at={(0.92,0.92)}, anchor=north east, font=\small},
  legend cell align=left,
]
\addplot[curveblue, very thick, domain=1:10, samples=250] {sin(deg(x))/x^2};
\addlegendentry{$y = \frac{\sin x}{x^2}$}
\addplot[gray!60, thick, dashed, domain=1:10, samples=100] {1/x^2};
\addlegendentry{envelope $\pm \frac{1}{x^2}$}
\addplot[gray!60, thick, dashed, domain=1:10, samples=100, forget plot] {-1/x^2};
\end{axis}
\end{tikzpicture}
\caption{Absolute convergence via an envelope. The oscillating curve $y = \frac{\sin x}{x^2}$ (blue) is squeezed between the symmetric boundaries $\pm \frac{1}{x^2}$ (dashed gray). Because the envelope encloses a finite total area ($\int_1^\infty \dd x/x^2 = 1$), the oscillations have no room to accumulate an infinite sum, guaranteeing that $\int_1^\infty \frac{\sin x}{x^2}\,\dd x$ converges absolutely.}
\label{ch25:fig-envelope}
\end{figure}

\section{A gallery of fully worked examples}
\label{ch25:gallery}

Here are a handful of complete computations, each touching a different technique.

\begin{example}[Integration by parts to infinity]
\label{ch25:ex-parts}
Evaluate $\displaystyle\int_0^{\infty} x\,e^{-x}\,\dd x$. Go to the wall $t$ and
integrate by parts ($u = x$, $\dd v = e^{-x}\dd x$, so $\dd u = \dd x$,
$v = -e^{-x}$):
\[
  \int_0^{t} x e^{-x}\,\dd x
  = \eval{-x e^{-x}}{0}^{t} + \int_0^{t} e^{-x}\,\dd x
  = -t e^{-t} + \bigl(1 - e^{-t}\bigr).
\]
As $t\to\infty$, the term $t e^{-t}\to 0$ (exponential decay beats the linear
growth of $t$ -- you proved this in the L'H\^opital chapter), and $e^{-t}\to 0$.
So the value is $0 + (1 - 0) = 1$. This is the mean of the exponential
distribution with rate $k=1$.
\end{example}

\begin{example}[A logarithmic spike]
\label{ch25:ex-logspike}
Evaluate $\displaystyle\int_0^{1} \ln x \, \dd x$. The integrand blows up at
$x=0^+$ (since $\ln x \to -\infty$), so this is second-kind. Back off and use parts
($u=\ln x$, $\dd v = \dd x$):
\[
  \int_{\eps}^{1}\ln x\,\dd x
  = \eval{x\ln x}{\eps}^{1} - \int_{\eps}^{1} 1\,\dd x
  = \bigl(1\cdot 0 - \eps\ln\eps\bigr) - (1 - \eps)
  = -\eps\ln\eps - 1 + \eps.
\]
As $\eps\to 0^+$, the product $\eps\ln\eps\to 0$ (again the L'H\^opital fact:
a power kills a logarithm), and $\eps\to 0$, leaving
\[
  \int_0^{1}\ln x\,\dd x = 0 - 1 + 0 = -1.
\]
A finite, negative area -- negative because $\ln x < 0$ on $(0,1)$ -- under a curve
that plunges to $-\infty$ at the left edge.
\end{example}

\begin{example}[Doubly improper: both kinds at once]
\label{ch25:ex-double-improper}
Evaluate $\displaystyle\int_0^{\infty}\frac{\dd x}{\sqrt{x}\,(1+x)}$. This is
improper \emph{twice}: the integrand blows up at $x=0$ \emph{and} the interval runs
to $\infty$. The rule is to split off each bad place. Split at $x=1$:
\[
  \int_0^{\infty} = \int_0^{1} + \int_1^{\infty},
\]
and demand both converge. Substitute $x = u^2$, $\dd x = 2u\,\dd u$, so
$\sqrt x = u$:
\[
  \int \frac{\dd x}{\sqrt x\,(1+x)}
  = \int \frac{2u\,\dd u}{u\,(1+u^2)}
  = 2\int \frac{\dd u}{1+u^2}
  = 2\arctan u = 2\arctan\sqrt x.
\]
Now both limits are clean. Near $0$: $2\arctan\sqrt x \to 0$. Near $\infty$:
$2\arctan\sqrt x \to 2\cdot \tfrac{\pi}{2} = \pi$. Each piece converges, so
\[
  \int_0^{\infty}\frac{\dd x}{\sqrt x\,(1+x)} = \pi - 0 = \pi.
\]
The substitution \emph{absorbed} both singularities at once, turning a doubly
improper integral into a single tidy arctangent.
\end{example}

\begin{whybox}
Why was it legitimate to use the antiderivative $2\arctan\sqrt x$ and just plug in
the endpoints $0$ and $\infty$, instead of writing the $\eps$ and $t$ limits
explicitly? Because once you have a single antiderivative that is \emph{continuous
all the way down to (and out to) the bad points} -- and $2\arctan\sqrt x$ is, it has
finite one-sided limits $0$ and $\pi$ -- the limit operations reduce to evaluating
those limits of the antiderivative. The $\eps$ and $t$ are still there; you have
just done them in your head. The danger, as the $\int_{-1}^1 \dd x/x^2$ trap showed,
is using an antiderivative that is \emph{not} continuous across the singularity.
When the antiderivative behaves nicely at the boundary, the shortcut is safe; when
it blows up there, you must keep the limits explicit and check convergence.
\end{whybox}

\section{What we built}
\label{ch25:summary}

We extended the integral past its original home -- bounded functions on bounded
intervals -- by the single device of \emph{taking a limit of ordinary integrals}.

\begin{itemize}
  \item \textbf{First kind} (infinite interval):
  $\int_a^\infty f = \lim_{t\to\infty}\int_a^t f$. Split doubly-infinite integrals
  into independent halves; never use a symmetric straddling limit.
  \item \textbf{Second kind} (unbounded integrand): back off from the singularity
  by $\eps$ and let $\eps\to 0^+$. Watch for \emph{interior} singularities that
  hide inside an innocent-looking interval.
  \item \textbf{The $p$-tests}, the two yardsticks: $\int_1^\infty \dd x/x^p$
  converges iff $p>1$; $\int_0^1 \dd x/x^p$ converges iff $p<1$. The inequality
  flips because the two ends reward opposite behavior.
  \item \textbf{Comparison and limit comparison} settle convergence even with no
  antiderivative: trap your curve under a finite ceiling or over an infinite floor;
  or match its leading behavior to a known $p$-integral.
  \item \textbf{Absolute convergence} handles sign-changing integrands: if
  $\int\abs{f}$ converges (checkable by comparison), so does $\int f$.
  \item \textbf{The surprises} -- finite area under an infinite curve, finite
  volume inside Gabriel's Horn's infinite surface -- are not paradoxes but lessons:
  area, length, volume, and surface are genuinely independent measurements, and the
  infinite must be \emph{computed}, never guessed.
\end{itemize}

The deepest takeaway is a habit of mind: an improper integral is never evaluated
``at infinity'' or ``at the spike.'' It is always an ordinary finite integral with a
limit wrapped around it. Keep the two steps separate, ask whether the tail is
thinner or fatter than $1/x$, and the entire subject organizes itself around that
one borderline curve.

\section{Applications and the road ahead}

\begin{connectionbox}
Improper integrals extend integration to unbounded intervals and unbounded functions,
making them essential across physics, probability, and analysis:
\begin{itemize}
  \item \emph{Probability:} Continuous probability distributions over unbounded domains,
  such as the normal distribution $\int_{-\infty}^{\infty} \frac{1}{\sqrt{2\pi}\sigma} e^{-(x-\mu)^2/(2\sigma^2)}\,\dd x = 1$,
  are defined by improper integrals.
  \item \emph{Physics:} Gravitational potential energy $-\int_\infty^r \frac{GMm}{s^2}\,\dd s$,
  electrostatic potentials, and escape velocity calculations integrate force over infinite paths.
  \item \emph{Integral transforms:} The Laplace and Fourier transforms integrate functions against
  decaying exponentials or sinusoids over unbounded domains.
\end{itemize}

In the next chapter, we apply integration to geometry and physics---computing volumes, arc
lengths, surface areas, and work. Later chapters develop numerical integration for integrands
without elementary antiderivatives (Chapter~27) and explore the discrete analogue of improper
integrals in infinite series (Chapter~28).
\end{connectionbox}

\chapterline

\section{Exercises}
\label{ch25:exercises}

\begin{exercise}[Direct evaluation, first kind]
\label{ch25:exr-eval-first}
Evaluate each by integrating to a finite wall and taking the limit, or show it
diverges:
\[
  \text{(a) } \int_2^{\infty}\frac{\dd x}{x^3},
  \qquad
  \text{(b) } \int_0^{\infty} e^{-3x}\,\dd x,
  \qquad
  \text{(c) } \int_1^{\infty}\frac{\dd x}{x^{4/5}}.
\]
\end{exercise}

\begin{exercise}[Direct evaluation, second kind]
\label{ch25:exr-eval-second}
Evaluate or show divergence, being explicit about where the integrand blows up:
\[
  \text{(a) } \int_0^{4}\frac{\dd x}{\sqrt{x}},
  \qquad
  \text{(b) } \int_0^{1}\frac{\dd x}{x^{3/2}},
  \qquad
  \text{(c) } \int_0^{2}\frac{\dd x}{(x-2)^{2/3}}.
\]
\end{exercise}

\begin{exercise}[The hidden interior singularity]
\label{ch25:exr-hidden}
A student computes $\int_{-2}^{2}\frac{\dd x}{x^4}$ by writing
$\eval{-\frac{1}{3x^3}}{-2}^{2}$ and reports a finite number. Explain in one or two
sentences why the answer is wrong, identify the singularity, and determine the
correct status (converge to what, or diverge?).
\end{exercise}

\begin{exercise}[Comparison test]
\label{ch25:exr-comparison}
Use the (plain) comparison test to decide convergence. State clearly which known
integral you compare to and in which direction.
\[
  \text{(a) } \int_1^{\infty}\frac{\dd x}{x^3 + 5},
  \qquad
  \text{(b) } \int_1^{\infty}\frac{2 + \cos x}{x}\,\dd x,
  \qquad
  \text{(c) } \int_1^{\infty}\frac{\dd x}{\sqrt{x^4 + 1}}.
\]
\end{exercise}

\begin{exercise}[Limit comparison]
\label{ch25:exr-limit-comparison}
Use limit comparison against a suitable $1/x^p$ to decide convergence:
\[
  \text{(a) } \int_1^{\infty}\frac{x + 1}{x^3 - x + 2}\,\dd x,
  \qquad
  \text{(b) } \int_1^{\infty}\frac{\dd x}{\sqrt{x^2 + 3x}}.
\]
\end{exercise}

\begin{exercise}[Absolute convergence]
\label{ch25:exr-absolute}
Show that $\displaystyle\int_1^{\infty}\frac{\cos x}{x^2 + x}\,\dd x$ converges by
proving it converges absolutely. (Bound $\abs{\cos x}$ and compare.)
\end{exercise}

\begin{exercise}[The $p$-spectrum at both ends]
\label{ch25:exr-both-ends}
Explain, using both $p$-tests, why $\displaystyle\int_0^{\infty}\frac{\dd x}{x^p}$
diverges for \emph{every} real $p$. (Split at $x=1$ and analyze each piece; show no
single $p$ saves both.)
\end{exercise}

\begin{exercise}[Gabriel's Horn, your turn]
\label{ch25:exr-horn}
Repeat the Gabriel's-Horn computation for the curve $y = 1/x^2$, $x\ge 1$, revolved
about the $x$-axis.
(a) Show the volume $\int_1^\infty \pi (1/x^2)^2\,\dd x$ is finite and find it.
(b) Explain why the lower bound $\sqrt{1+(y')^2}\ge 1$ fails to decide convergence, and decide whether the surface area
$\int_1^\infty 2\pi\,\frac{1}{x^2}\sqrt{1+(y')^2}\,\dd x$ converges or diverges this
time. (Hint: the square root $\sqrt{1+(2/x^3)^2}$ is bounded for $x\ge 1$, so
compare to $\int_1^\infty 2\pi/x^2$.) Contrast with the $y=1/x$ horn: why does this
one \emph{not} give the paradox?
\end{exercise}

\begin{exercise}[Proving the Limit Comparison Test]
\label{ch25:exr-prove-lct}
Suppose $f,g>0$ on the interval $[a,\infty)$ and $\lim_{x\to\infty} f(x)/g(x) = L$ with
$0<L<\infty$. Using the definition of the limit, show that there is some $X$ beyond
which $\frac{L}{2}\,g(x) \le f(x) \le 2L\,g(x)$. Then invoke the plain comparison
test (twice) to conclude $\int f$ and $\int g$ share the same fate. (This is a
clean, complete proof -- write out both directions.)
\end{exercise}

\begin{exercise}[A convergence that needs L'H\^opital]
\label{ch25:exr-lhopital}
Evaluate $\displaystyle\int_1^{\infty}\frac{\ln x}{x^2}\,\dd x$. Integrate by parts
to the wall $t$, then take $t\to\infty$, using $\lim_{t\to\infty}\frac{\ln t}{t}=0$.
(Answer: $1$.)
\end{exercise}

\begin{exercise}[Conceptual: why split?]
\label{ch25:exr-why-split}
In two or three sentences, explain why $\int_{-\infty}^{\infty} x\,\dd x$ is said to
diverge even though $\lim_{t\to\infty}\int_{-t}^{t} x\,\dd x = 0$. What is the
difference between the principal value and genuine convergence, and why do we
insist on the stricter notion?
\end{exercise}

\begin{exercise}[A surprise of your own making]
\label{ch25:exr-surprise}
Find a continuous function $f$ on $[1,\infty)$ that does \emph{not} tend to zero
as $x\to\infty$, yet has $\int_1^\infty f\,\dd x$ convergent.
(Hint: tall thin spikes. Put a triangular spike of height $1$ and base width
$1/2^n$ centered at $x=n$. The integrand keeps hitting $1$ -- so it fails to go to
zero pointwise everywhere -- yet the total area is a convergent geometric series.
This shows the ``$f\to0$'' folklore is not even necessary in the strict pointwise
sense; only the \emph{tail area} must vanish.)
\end{exercise}
